% Figure from MATH 590 notes, section 22. Compile with the notes preamble.
\begin{tikzpicture}[scale=1.25,
  declare function={fx(\u)=3*\u; fy(\u)=0.3*sin(360*\u)+0.2;
                    gx(\u)=0.5+2.6*\u; gy(\u)=2.1+0.3*cos(540*\u);}]
  % trajectories t -> H(x,t)
  \foreach \u in {0,0.1,...,1.01} {
    \draw[black!30, thin] ({fx(\u)},{fy(\u)}) -- ({gx(\u)},{gy(\u)});
  }
  % intermediate maps H(., 1/3), H(., 2/3)
  \foreach \t in {0.333,0.667} {
    \draw[black!50, dashed, variable=\u, domain=0:1, samples=60, smooth]
      plot ({(1-\t)*fx(\u)+\t*gx(\u)}, {(1-\t)*fy(\u)+\t*gy(\u)});
  }
  % f and g
  \draw[figB, thick, variable=\u, domain=0:1, samples=60, smooth] plot ({fx(\u)},{fy(\u)});
  \draw[figB, thick, variable=\u, domain=0:1, samples=60, smooth] plot ({gx(\u)},{gy(\u)});
  \node[font=\small, figB, anchor=west] at ({fx(1)+0.08},{fy(1)}) {$f$};
  \node[font=\small, figB, anchor=west] at ({gx(1)+0.08},{gy(1)}) {$g$};
  \node[font=\scriptsize, black!60, anchor=west] at ({(2*fx(1)+gx(1))/3+0.08},{(2*fy(1)+gy(1))/3}) {$H(\cdot, \tfrac13)$};
  \node[font=\scriptsize, black!60, anchor=west] at ({(fx(1)+2*gx(1))/3+0.08},{(fy(1)+2*gy(1))/3}) {$H(\cdot, \tfrac23)$};
  % one highlighted trajectory
  \def\u{0.35}
  \draw[figR, thick] ({fx(\u)},{fy(\u)}) -- ({gx(\u)},{gy(\u)});
  \fill[figR] ({fx(\u)},{fy(\u)}) circle (1.4pt) node[below left, font=\scriptsize, inner sep=1pt, fill=white, xshift=-2pt, yshift=-2pt] {$f(x)$};
  \fill[figR] ({gx(\u)},{gy(\u)}) circle (1.4pt) node[above, font=\scriptsize, inner sep=1pt, yshift=3pt] {$g(x)$};
  \fill[figR] ({0.5*fx(\u)+0.5*gx(\u)},{0.5*fy(\u)+0.5*gy(\u)}) circle (1.4pt)
    node[left, font=\scriptsize, fill=white, inner sep=1pt, xshift=-3pt] {$H(x,\tfrac12)$};
  \node[font=\small] at (-0.35,2.5) {$\mathbb{R}^n$};
\end{tikzpicture}
